% ============================================================
%  II. PROBLEM STATEMENT
% ============================================================
\section{Problem Statement}
\label{sec:problem_formulation}

\subsection{System Model}

The \ac{TERESA} platform consists of a Boston Dynamics Spot quadruped and a
Unitree Z1 six-\ac{DOF} arm rigidly mounted on the robot's back. Four frames
are defined: the world frame $\mathcal{F}_W$ (coincident with Spot odometry),
the body frame $\mathcal{F}_B$, the arm base frame $\mathcal{F}_0$ (fixed
relative to $\mathcal{F}_B$), and the end-effector frame $\mathcal{F}_{EE}$ (at
the Z1 flange). The Spot base is modelled as a non-holonomic vehicle with
configuration $\boldsymbol{\xi} = [x,\, y,\, \theta]^\top \in \mathbb{R}^2
\times \mathbb{S}^1$, constrained to forward velocity $v_x$ and yaw rate
$\omega_z$:
\begin{equation}
\dot{\boldsymbol{\xi}} =
\begin{bmatrix} \cos\theta \\ \sin\theta \\ 0 \end{bmatrix} v_x
+
\begin{bmatrix} 0 \\ 0 \\ 1 \end{bmatrix} \omega_z.
\label{eq:nonholonomic}
\end{equation}

The body posture is parametrised by height $h$ and pitch $\varphi$, commanded
through the native posture control interface, with bounds
$h \in [h_{\min},\, h_{\max}]$, $\varphi \in [0,\, \varphi_{\max}]$. The Z1 arm
has $n = 6$ revolute joints $\mathbf{q} \in \mathbb{R}^6$ with geometric
Jacobian $\mathbf{J}_{\mathrm{arm}}(\mathbf{q}) \in \mathbb{R}^{6 \times 6}$.
The combined control input during the approach phase is
$\mathbf{u} = [\dot{\mathbf{q}},\, v_x,\, \omega_z]^\top \in \mathbb{R}^8$.

\subsection{Task Definition}

An autonomous emergency assessment mission requires the \ac{EE} to reach a
sequence of $N$ anatomical target poses
$\mathcal{T} = \{\mathbf{T}_i^* \in SE(3)\}_{i=1}^{N}$ (e.g., \ac{FAST}
quadrants). The targets are not known a priori---they are estimated during the
mission from the onboard \ac{RGBD} cameras, which are the only source of
perception.

\subsection{Problem}

Given the system model and the unknown patient position, find control inputs
$\mathbf{u}(t)$ and body posture commands $(h(t), \varphi(t))$ that enable the
robot to: (i)~discover a posture from which detection confidence exceeds a
threshold, ensuring robust patient localisation before committing to the
approach; (ii)~navigate to a configuration from which all targets in
$\mathcal{T}$ are reachable; (iii)~reconfigure its body to maximise arm
manipulability for each scanning target; and (iv)~bring the \ac{EE} to each
target pose with bounded error $\|\mathbf{p}_{EE} - \mathbf{p}_i^*\| \leq
\varepsilon_p$.

\subsection{Assumptions}

The following assumptions are adopted: the patient remains stationary during
scanning (A1); Spot operates on flat terrain (A2); the rigid transform between
$\mathcal{F}_B$ and $\mathcal{F}_0$ is known and constant (A3); a single
patient is present (A4); and no external perception infrastructure is available
(A5).
